\documentclass[10pt, a4paper]{article}

\usepackage[utf8]{inputenc}
\usepackage[T1]{fontenc}
\usepackage[french]{babel}
\usepackage[left=1cm, right=1cm, top=1cm, bottom=1cm]{geometry}
\usepackage{titlesec}
\usepackage{enumitem}
\usepackage{xcolor}
\usepackage{fontawesome5}
\usepackage{hyperref}

% Couleurs personnalisées
\definecolor{primary}{RGB}{0, 64, 128} % Bleu foncé élégant pour les titres
\definecolor{textgray}{RGB}{50, 50, 50}

% Configuration des liens hypertexes (colore le texte sans encadré)
\hypersetup{
    colorlinks=true,
    urlcolor=primary,
    linkcolor=primary,
    pdfborder={0 0 0}
}

% Configuration compacte des sections et des listes
\titleformat{\section}{\large\bfseries\color{primary}}{}{0em}{}[\titlerule]
\titlespacing*{\section}{0pt}{1.2ex}{0.8ex}
\setlist[itemize]{
    label=\small\textbullet,
    leftmargin=1.5em,
    itemsep=0.2ex,
    parsep=0pt,
    topsep=0.2ex,
    partopsep=0pt
}

% En-têtes d'expériences et de formations
\newcommand{\cventry}[4]{%
    \noindent\textbf{#1} \hfill \textit{#2} \\
    \textit{#3} \hfill {\footnotesize #4} \vspace{0.2ex}
}

\newcommand{\experience}[3]{%
    \noindent
    \parbox[t]{0.72\textwidth}{\textbf{#1} \textit{\,|\,#2}}%
    \hfill
    \parbox[t]{0.26\textwidth}{\raggedleft\normalfont\itshape\mbox{#3}}%
    \par\vspace{0.2ex}
}

\begin{document}
\pagestyle{empty} % Désactive les numéros de page

% ------------------- EN-TÊTE -------------------
\begin{center}
	{\huge \textbf{Tsikomia Scidelo ANDRIASAHY}}\\[0.5ex]
	{\large \textcolor{primary}{\textbf{Élève-ingénieur en Systèmes Électroniques et Automatiques}}}\\[0.5ex]
	\small
	\faMapMarker \ Tanger, Maroc (Mobilité nationale \& internationale --- Stage PFE) \quad
	\faPhone \ +212 777 647 362 \quad
	\faEnvelope \ \href{mailto:scideloandriasahy@gmail.com}{scideloandriasahy@gmail.com} \\[0.5ex]
	\faLinkedin \ \href{https://linkedin.com/in/scidelo}{linkedin.com/in/scidelo} \quad
	\faGithub \ \href{https://scidelo-andriasahy.github.io/}{scidelo-andriasahy.github.io}
\end{center}

\vspace{-1ex}

% ------------------- PROFIL -------------------
\section*{Profil}
\noindent Élève-ingénieur GSEA à l'\href{https://ensat.ac.ma/Portail/}{\textbf{ENSA Tanger}}, spécialisé en instrumentation, traitement du signal et \mbox{automatique}. Fort de réalisations concrètes en R\&D (bancs de mesure LabVIEW, diagnostic par imagerie thermique/IA et régulation avancée), j'associe maîtrise des capteurs physiques et analyse prédictive des données sous Python et Simulink.

\noindent Disponible dès \textbf{mi-janvier 2027} pour un stage de fin d'études (PFE), ouvert à toute modalité et mobile au Maroc comme à l'international, afin de développer des outils intelligents d'aide à la décision et fiabiliser des systèmes complexes.

% ------------------- FORMATION -------------------
\section*{Formation}

\cventry{\href{https://ensat.ac.ma/Portail/genie-des-systemes-electroniques-et-automatiques-gsea/}{Cycle d'Ingénieur – Génie des Systèmes Électroniques et Automatiques (GSEA)}}{2022 -- 2027}
{\href{https://www.uae.ac.ma/}{Université Abdelmalek Essaâdi} (\href{https://ensat.ac.ma/Portail/}{ENSA}), Tanger -- \textit{3\textsuperscript{ème} année (Dernière année)}}{}
\begin{itemize}
	\item \textbf{Spécialités :} Systèmes embarqués, électronique de puissance, microcontrôleurs, traitement du signal, automatique.
\end{itemize}
\vspace{0.6ex}

\cventry{\href{https://facsciences.univ-fianarantsoa.mg/fr/formations/parcours/vorhf2gwn3pfzqta7w2pff6p}{Licence 1 -- Mathématiques et Informatique pour les Sciences Sociales (MISS)}}{2021 -- 2022}
{\href{https://www.univ-fianarantsoa.mg/}{Université de Fianarantsoa}, Madagascar}{}
\vspace{0.6ex}

\cventry{\textcolor{primary}{Baccalauréat Scientifique\ \textbf{---}\ Série C \textnormal{(Mention Très Bien)}}}{2020 -- 2021}
{\href{https://www.facebook.com/p/Coll\%C3\%A8ge-Priv\%C3\%A9-Houssen-Memorial-School-61557463046217/}{Houssen Memorial School (HMS)}, Tuléar, Madagascar}{}

% ------------------- COMPÉTENCES TECHNIQUES -------------------
\section*{Compétences Techniques}
\begin{itemize}[leftmargin=*, itemsep=0.2ex, parsep=0pt]
	\item \textbf{Systèmes Embarqués \& Électronique :} STM32, PIC, ESP32, Arduino, DSP, capteurs industriels, KiCad, Proteus.
	\item \textbf{Programmation \& Firmware :} C (MISRA C), C\texttt{++}, Python, VHDL, développement de Firmware embarqué, PlatformIO.
	\item \textbf{Automatique \& IA Embarquée :} Régulation PID, Logique floue, ANFIS, Algorithmes Génétiques (GA), Commande par réseau de neurones (ANN), Algorithme MPPT P\&O, CNN, Modèles de Markov.
	\item \textbf{Outils :} MATLAB/Simulink, ROS2, TIA Portal, LaTeX, LabVIEW, Linux, Git/GitHub, Android Studio.
\end{itemize}

% ------------------- EXPÉRIENCE PROFESSIONNELLE -------------------
\section*{Expérience Professionnelle}

\experience{Ingénieur Stagiaire R\&D Électronique de Puissance}{\href{https://sitewebfortico.vercel.app/fr}{Fortico}}{Juil. 2026 - Oct. 2026}
\begin{itemize}
	\item Conception (KiCad/Proteus) d'un onduleur hybride : MPPT entrelacé 3 phases et étages 24V/350VDC/230VAC.
	\item Simulation Simulink et régulation PID : MPPT P\&O, profil de charge batterie et gestion hybride réseau.
\end{itemize}

\experience{Stagiaire R\&D IA \& IoT (PFA)}{\href{https://www.linkedin.com/company/smart-automation-technologies/home/}{Smart Automation Technologies}}{Juil. 2026 - Août 2026}
\begin{itemize}
	\item Analyse comparative de modèles d'IA (Markov vs CNN) pour le diagnostic par imagerie thermique.
	\item Déploiement sur nœuds IoT pour la détection automatique d'anomalies et la maintenance prédictive.
\end{itemize}

\experience{Stage d'observation en Électricité}{\href{https://www.amendis.ma/}{Amendis} (Groupe Veolia)}{Juil. 2025 - Août 2025}
\begin{itemize}
	\item Découverte de la supervision SCADA et des protocoles de communication pour l'automatisation.
\end{itemize}

\experience{Animateur d'Ateliers Robotique}{\href{https://www.facebook.com/steamacademy.ma/}{STEAM Academy}}{2025 - Présent}
\begin{itemize}
	\item Encadrement d'enfants autour des concepts de base en robotique et de la programmation.
\end{itemize}

\experience{Professeur Particulier (Sciences \& Technologie)}{\href{https://www.superprof.ma/}{Superprof} / Indépendant}{2024 - Présent}

\experience{Rédacteur Technique (Freelance)}{\href{https://www.alphorm.com/}{Alphorm}}{Août 2024 - Mars 2025}

% ------------------- PROJETS TECHNIQUES & ACADÉMIQUES -------------------
\section*{Projets Techniques \& Académiques}
\begin{itemize}[leftmargin=*, itemsep=0.4ex, parsep=0pt]
	\item {\color{primary}\textbf{Instrumentation Virtuelle \& Oscilloscope Numérique (LabVIEW) :}}
	      Développement d'un banc de mesure temps réel (générateur multi-signaux, gestion d'échelles Volt/Div, ms/Div).

	\item {\color{primary}\textbf{Robot Mobile Autonome AMR (ROS 2 Humble, Ubuntu) :}}
	      Conception du jumeau numérique sous Gazebo (URDF/Xacro) avec simulation capteurs (LiDAR 2D, IMU).
	      Implémentation du Graph-SLAM (SLAM Toolbox, Ceres) et navigation autonome via la stack Nav2 (Costmaps,
	      planification globale/locale A*/DWB), validée par déploiement physique et synchronisation des repères TF.

	\item {\color{primary}\textbf{Commande \& Suivi de Trajectoire d'un Bras Manipulateur (2 DDL) :}}
	      Modélisation dynamique sous MATLAB/Simulink et accélération par S-Functions C-MEX ($\times$50).
	      Étude comparative entre réseau de neurones (ANN), PID optimisé par algorithme génétique (GA) et logique
	      floue (FLC) : sélection de la commande floue, validation logicielle en C++ et déploiement sur carte DSP TMS320C6713.
\end{itemize}

% ------------------- CERTIFICATIONS & EXPÉRIENCE ASSOCIATIVE -------------------
\section*{Certifications, Expérience Associative, Langues \& Divers}
\begin{itemize}[leftmargin=*, itemsep=0.3ex, parsep=0pt]
	\item \textbf{Certifications \& Récompenses :} Formation certifiante en IA \& Idéation entrepreneuriale (\href{http://cidealmaroc.org/}{Fondation CIDEAL} -- \mbox{Programme} INNO'VERT / ENSA Tanger) \textbar{} Formation certifiante validée par \href{https://web.facebook.com/profile.php?id=61567878936369}{IngénieriePro} \textbar{} Participation certifiée à l'\href{https://iymc.info/fr/}{IYMC} (International Youth Math Challenge) \textbar{} Participation officielle au \href{https://www.fitut.ma/}{FITUT} (Théâtre universitaire en 2022).
	\item \textbf{Associatif :} Adjoint chargé académique à l'\href{https://www.linkedin.com/company/eisclubuae/home/}{\textbf{EISCLUB}} (2023 -- 2024). Responsable Logistique pour \href{https://tsingymaroc.org/}{\textbf{TSINGY TANGER}} (2022 -- 2024). Membre actif : \href{https://cesamcentrale.org/category/sections-cesam/tanger/}{CESAM Tanger}, \href{https://www.instagram.com/ensatronics/}{EnsaTronics} (Robotique) et \href{https://www.ojafricains.org/}{OJA} (Organisation des Jeunes Africains).
	\item \textbf{Langues :} Français (Courant), Anglais (Technique, lu et écrit), Espagnol (Notions), Malagasy (Langue maternelle).
	\item \textbf{Centres d'intérêt :} Cinéphilie, veille technologique (salons \& conférences), jeu d'échecs, pétanque, football, DIY.
\end{itemize}
\end{document}
